\scanpage{58}
\mspage{4}

puisse avoir \deleted{\unsure{aussi}} même
\[
\tau_1r=\tau_\infty r=\tau_0r
\]
i.e. qu'on soit dans le cas (10) (purement comb.: gr. d'op. $G_2$, avec les $\unsure{\sigma_i|\tau_i}$ opérant sans pts fixes) cela signifie que le point de \uline{type 1} du repère envisagé est un pli (les points de type 0 ne peuvent l'être, par définition).

Grâce à l'introduction des arêtes pliées, dans la équiv. de catégories (8), il n'est plus néc. d'introduire l'hyp que $\rho_1$ opère sans pt fixe -- et dans l'équivalence (1), il n'y a plus lieu de supposer que $\rho_1$ (alias $\sigma$) opère ss pt fixes -- par contre la condition que les $\unsure{\sigma_i|\tau_i}$ opèrent sans points fixes demeure provisoirement (mais il serait facile maintenant de s'en affranchir \deleted{\unsure{?}} cf plus bas\ldots)).

\textcircled{2} On a dans ces considérations un rôle particulier pour l'indice $1\in\{0,1,\infty\}$, du fait que l'on a
\[
\rho_1^2=1,
\]
on peut s'en défaire, en réintroduisant
\[
(12)\qquad\Pi_{03}^{\deleted{\unsure{\omega}}\tau}=\{\tau_0,\tau_1,\tau_\infty\mid\tau_0^2=\tau_1^2=\tau_\infty^2=1\}
\]
(sans condition $(\tau_0\tau_\infty)^2=1$ !), et le sous-groupe d'indice 2
\margin{$\Pi_{03}$ \unsure{est le noyau de} \deleted{$\unsure{\Pi_{03}}$}
\diagram{\begin{tikzcd}[column sep=small,row sep=small]
\Pi_{03}^{\tau}\arrow[r]&\mathbb Z/2\\
\unsure{\sigma_i|\tau_i}\arrow[r,mapsto]&1\bmod2
\end{tikzcd}}}
\[
(13)\qquad\Pi_{03}=\{\rho_0,\rho_1,\rho_\infty\mid\rho_\infty\rho_1\rho_0=1\}
\]
\mspage{57}
